\documentclass[10pt,a4paper]{amsart}
\usepackage{amsmath,amssymb,amsthm,graphicx}
% Retained minimal preamble for the Stacks Project chapter flat.tex (More on Flatness).
% Verbatim excerpt of the chapter preamble: Xy-pic setup (native preamble lines 14-18) and the native
% theorem declarations with the mathematical macros (native preamble lines 154-217). Copied mechanically;
% nothing added, removed or redefined. The wrapper class and the amsmath/amssymb/amsthm/graphicx base
% packages are supplied by the assembler, so the class/externaldocument part of the native preamble is omitted.
% For commutative diagrams we use Xy-pic
\usepackage[all]{xy}

% We use 2cell for 2-commutative diagrams.
\xyoption{2cell}
% Theorem environments.
%
\theoremstyle{plain}
\newtheorem{theorem}[subsection]{Theorem}
\newtheorem{proposition}[subsection]{Proposition}
\newtheorem{lemma}[subsection]{Lemma}

\theoremstyle{definition}
\newtheorem{definition}[subsection]{Definition}
\newtheorem{example}[subsection]{Example}
\newtheorem{exercise}[subsection]{Exercise}
\newtheorem{situation}[subsection]{Situation}

\theoremstyle{remark}
\newtheorem{remark}[subsection]{Remark}
\newtheorem{remarks}[subsection]{Remarks}

\numberwithin{equation}{subsection}

% Macros
%
\def\lim{\mathop{\mathrm{lim}}\nolimits}
\def\colim{\mathop{\mathrm{colim}}\nolimits}
\def\Spec{\mathop{\mathrm{Spec}}}
\def\Hom{\mathop{\mathrm{Hom}}\nolimits}
\def\Ext{\mathop{\mathrm{Ext}}\nolimits}
\def\SheafHom{\mathop{\mathcal{H}\!\mathit{om}}\nolimits}
\def\SheafExt{\mathop{\mathcal{E}\!\mathit{xt}}\nolimits}
\def\Sch{\mathit{Sch}}
\def\Mor{\mathop{\mathrm{Mor}}\nolimits}
\def\Ob{\mathop{\mathrm{Ob}}\nolimits}
\def\Sh{\mathop{\mathit{Sh}}\nolimits}
\def\NL{\mathop{N\!L}\nolimits}
\def\CH{\mathop{\mathrm{CH}}\nolimits}
\def\proetale{{pro\text{-}\acute{e}tale}}
\def\etale{{\acute{e}tale}}
\def\QCoh{\mathit{QCoh}}
\def\Ker{\mathop{\mathrm{Ker}}}
\def\Im{\mathop{\mathrm{Im}}}
\def\Coker{\mathop{\mathrm{Coker}}}
\def\Coim{\mathop{\mathrm{Coim}}}

% Boxtimes
%
\DeclareMathSymbol{\boxtimes}{\mathbin}{AMSa}{"02}

%
% Macros for moduli stacks/spaces
%
\def\QCohstack{\mathcal{QC}\!\mathit{oh}}
\def\Cohstack{\mathcal{C}\!\mathit{oh}}
\def\Spacesstack{\mathcal{S}\!\mathit{paces}}
\def\Quotfunctor{\mathrm{Quot}}
\def\Hilbfunctor{\mathrm{Hilb}}
\def\Curvesstack{\mathcal{C}\!\mathit{urves}}
\def\Polarizedstack{\mathcal{P}\!\mathit{olarized}}
\def\Complexesstack{\mathcal{C}\!\mathit{omplexes}}
% \Pic is the operator that assigns to X its picard group, usage \Pic(X)
% \Picardstack_{X/B} denotes the Picard stack of X over B
% \Picardfunctor_{X/B} denotes the Picard functor of X over B
\def\Pic{\mathop{\mathrm{Pic}}\nolimits}
\def\Picardstack{\mathcal{P}\!\mathit{ic}}
\def\Picardfunctor{\mathrm{Pic}}
\def\Deformationcategory{\mathcal{D}\!\mathit{ef}}

% --- appended 2026-10-12 (Chief): the project own definitions of the reference/slogan/history
% display environments, copied mechanically from the shared Stacks preamble lines 8-12
% (\usepackage{verbatim} plus the three \newenvironment lines, which the chapter preamble keeps
% verbatim from the project preamble). Nothing added, removed or redefined: these environments
% carry no mathematical content and the project itself typesets them as comments. Without these
% three lines a counted statement whose own text uses them fails with Environment slogan undefined.
\usepackage{verbatim}
\newenvironment{reference}{\comment}{\endcomment}
\newenvironment{slogan}{\comment}{\endcomment}
\newenvironment{history}{\comment}{\endcomment}
\usepackage{cleveref}
\crefname{originalreference}{original reference}{original references}
\Crefname{originalreference}{Original reference}{Original references}
\makeatletter
\newlabel{situation-complex-and-divisor}{{0F8S}{source}{}{sourcepointer-1}{}}
\newlabel{situation-complex-and-divisor@cref}{{[originalreference][0][]0F8S}{[1][0][]source}{}{}{}}
\makeatother
\crefname{originalsection}{original section}{original sections}
\Crefname{originalsection}{Original section}{Original sections}
\makeatletter
\newlabel{more-algebra-section-perfect-eta}{{0F7V}{source}{}{sourcepointer-2}{}}
\newlabel{more-algebra-section-perfect-eta@cref}{{[originalsection][0][]0F7V}{[1][0][]source}{}{}{}}
\makeatother
\crefname{originallemma}{original lemma}{original lemmas}
\Crefname{originallemma}{Original lemma}{Original lemmas}
\makeatletter
\newlabel{divisors-lemma-blow-up-pullback-effective-Cartier}{{0809}{source}{}{sourcepointer-3}{}}
\newlabel{divisors-lemma-blow-up-pullback-effective-Cartier@cref}{{[originallemma][0][]0809}{[1][0][]source}{}{}{}}
\makeatother
\crefname{originallemma}{original lemma}{original lemmas}
\Crefname{originallemma}{Original lemma}{Original lemmas}
\makeatletter
\newlabel{divisors-lemma-blowing-up-two-ideals}{{080A}{source}{}{sourcepointer-4}{}}
\newlabel{divisors-lemma-blowing-up-two-ideals@cref}{{[originallemma][0][]080A}{[1][0][]source}{}{}{}}
\makeatother
\crefname{originallemma}{original lemma}{original lemmas}
\Crefname{originallemma}{Original lemma}{Original lemmas}
\makeatletter
\newlabel{divisors-lemma-blowing-up-gives-effective-Cartier-divisor}{{02OS}{source}{}{sourcepointer-5}{}}
\newlabel{divisors-lemma-blowing-up-gives-effective-Cartier-divisor@cref}{{[originallemma][0][]02OS}{[1][0][]source}{}{}{}}
\makeatother
\crefname{originallemma}{original lemma}{original lemmas}
\Crefname{originallemma}{Original lemma}{Original lemmas}
\makeatletter
\newlabel{lemma-pullback-triple-ideals-good}{{0F8T}{source}{}{sourcepointer-6}{}}
\newlabel{lemma-pullback-triple-ideals-good@cref}{{[originallemma][0][]0F8T}{[1][0][]source}{}{}{}}
\makeatother
\crefname{originallemma}{original lemma}{original lemmas}
\Crefname{originallemma}{Original lemma}{Original lemmas}
\makeatletter
\newlabel{more-algebra-lemma-eta-change-unit}{{0GST}{source}{}{sourcepointer-7}{}}
\newlabel{more-algebra-lemma-eta-change-unit@cref}{{[originallemma][0][]0GST}{[1][0][]source}{}{}{}}
\makeatother
\crefname{originallemma}{original lemma}{original lemmas}
\Crefname{originallemma}{Original lemma}{Original lemmas}
\makeatletter
\newlabel{more-algebra-lemma-ideal-well-defined}{{0GSS}{source}{}{sourcepointer-8}{}}
\newlabel{more-algebra-lemma-ideal-well-defined@cref}{{[originallemma][0][]0GSS}{[1][0][]source}{}{}{}}
\makeatother
\crefname{originallemma}{original lemma}{original lemmas}
\Crefname{originallemma}{Original lemma}{Original lemmas}
\makeatletter
\newlabel{divisors-lemma-universal-property-blowing-up}{{0806}{source}{}{sourcepointer-9}{}}
\newlabel{divisors-lemma-universal-property-blowing-up@cref}{{[originallemma][0][]0806}{[1][0][]source}{}{}{}}
\makeatother
\crefname{originallemma}{original lemma}{original lemmas}
\Crefname{originallemma}{Original lemma}{Original lemmas}
\makeatletter
\newlabel{constructions-lemma-relative-glueing}{{01LH}{source}{}{sourcepointer-10}{}}
\newlabel{constructions-lemma-relative-glueing@cref}{{[originallemma][0][]01LH}{[1][0][]source}{}{}{}}
\makeatother
\title{Existence and universal property of the blowup making a complex and divisor into a good triple}
\author{The Stacks Project Authors}
\date{stacks-project:flat.tex:local-snapshot-20261011; GNU FDL 1.2 or later}
\begin{document}
\maketitle
\noindent\textit{Editorial note.} All delivered mathematical text is the author's own, copied from the pinned Stacks Project chapter file flat.tex; no statement or proof step was rewritten, reordered or removed.\par\medskip
\noindent\textit{Editorial note.} The preamble is the chapter's own retained minimal preamble (sources/stacks-flat/retained-preamble.tex), a verbatim excerpt of the native Xy-pic setup and of the native theorem declarations and mathematical macros, followed by the project own definitions of the reference, slogan and history display environments (verbatim package plus three newenvironment lines) that the retained file carries at its end and that the project itself typesets as comments. No macro was invented or redefined, and the wrapper class and the amsmath/amssymb/amsthm/graphicx base packages are supplied by the assembler.\par\medskip
\noindent\textit{Editorial note.} The counted statement and proof contain no cross reference, no citation and no included file, so the delivered body is self contained.\par\medskip
\noindent\textit{Editorial note.} Canonical location of the counted proposition: Stacks tag 0GTY (sources/stacks-flat/pinned-original-source.tex line 13588).\par\medskip
\noindent\textit{Allowed standard prerequisites (editorial).}\begin{itemize}
\item Good triples and the ideals I\_i(M\textasciicircum{}bullet, f) attached to an affine chart. Same chapter, Situation labelled situation-complex-and-divisor and Lemma labelled lemma-pullback-triple-ideals-good, cited at proof lines 13632-13634; More on Algebra, Section labelled more-algebra-section-perfect-eta, cited at proof lines 13608-13610. Fixes the data and the ideals whose blowup is being constructed.
\item Blowing up in a product of ideals and the universal property of a blowup. Divisors, Lemma labelled divisors-lemma-blowing-up-two-ideals, cited at proof lines 13622-13624 and used for the factorization, and Divisors, Lemma labelled divisors-lemma-universal-property-blowing-up, cited at proof lines 13662. Factors the blowup through the individual blowups and gives the unique factorization of h.
\item Behaviour of the ideals under pullback and under a change of the equation of the divisor. More on Algebra, Lemmas labelled more-algebra-lemma-eta-change-unit, more-algebra-lemma-ideal-well-defined and lemma-pullback-triple-ideals-good, cited at proof lines 13652-13660 and 13700-13708; Divisors, Lemma labelled divisors-lemma-blow-up-pullback-effective-Cartier, cited at proof lines 13618-13620. Identifies the ideals computed in different charts and after pullback.
\item Gluing of relative schemes. Constructions, Lemma labelled constructions-lemma-relative-glueing, cited at proof lines 13680-13682. Descends the local morphisms b\_U to a morphism b : X' -> X.
\end{itemize}\medskip
\ifdefined\hypertarget\hypertarget{sourcepointer-1}{}\fi
\noindent\textit{Source reference.} Original reference 0F8S: sources/stacks-flat/pinned-original-source.tex, source lines 13460--13460. Same-chapter pointer into the author environment situation at source line 13460; invoked at proof line None of the counted proof. It is retained only as a navigation pointer; every load-bearing step is inside the delivered span.\par\medskip
\ifdefined\hypertarget\hypertarget{sourcepointer-2}{}\fi
\noindent\textit{Source reference.} Original section 0F7V: sources/stacks-chow/tags.txt, source lines 19020--19020. Cross-chapter standard prerequisite located in the project tag list staged in the local snapshot; invoked at proof line 13609 of the counted proof. It is retained only as a navigation pointer; every load-bearing step is inside the delivered span.\par\medskip
\ifdefined\hypertarget\hypertarget{sourcepointer-3}{}\fi
\noindent\textit{Source reference.} Original lemma 0809: sources/stacks-chow/tags.txt, source lines 10174--10174. Cross-chapter standard prerequisite located in the project tag list staged in the local snapshot; invoked at proof line 13627 of the counted proof. It is retained only as a navigation pointer; every load-bearing step is inside the delivered span.\par\medskip
\ifdefined\hypertarget\hypertarget{sourcepointer-4}{}\fi
\noindent\textit{Source reference.} Original lemma 080A: sources/stacks-chow/tags.txt, source lines 10175--10175. Cross-chapter standard prerequisite located in the project tag list staged in the local snapshot; invoked at proof line 13630 of the counted proof. It is retained only as a navigation pointer; every load-bearing step is inside the delivered span.\par\medskip
\ifdefined\hypertarget\hypertarget{sourcepointer-5}{}\fi
\noindent\textit{Source reference.} Original lemma 02OS: sources/stacks-chow/tags.txt, source lines 3527--3527. Cross-chapter standard prerequisite located in the project tag list staged in the local snapshot; invoked at proof line 13637 of the counted proof. It is retained only as a navigation pointer; every load-bearing step is inside the delivered span.\par\medskip
\ifdefined\hypertarget\hypertarget{sourcepointer-6}{}\fi
\noindent\textit{Source reference.} Original lemma 0F8T: sources/stacks-flat/pinned-original-source.tex, source lines 13488--13488. Same-chapter pointer into the author environment lemma at source line 13488; invoked at proof line 13639 of the counted proof. It is retained only as a navigation pointer; every load-bearing step is inside the delivered span.\par\medskip
\ifdefined\hypertarget\hypertarget{sourcepointer-7}{}\fi
\noindent\textit{Source reference.} Original lemma 0GST: sources/stacks-chow/tags.txt, source lines 20943--20943. Cross-chapter standard prerequisite located in the project tag list staged in the local snapshot; invoked at proof line 13650 of the counted proof. It is retained only as a navigation pointer; every load-bearing step is inside the delivered span.\par\medskip
\ifdefined\hypertarget\hypertarget{sourcepointer-8}{}\fi
\noindent\textit{Source reference.} Original lemma 0GSS: sources/stacks-chow/tags.txt, source lines 20942--20942. Cross-chapter standard prerequisite located in the project tag list staged in the local snapshot; invoked at proof line 13653 of the counted proof. It is retained only as a navigation pointer; every load-bearing step is inside the delivered span.\par\medskip
\ifdefined\hypertarget\hypertarget{sourcepointer-9}{}\fi
\noindent\textit{Source reference.} Original lemma 0806: sources/stacks-chow/tags.txt, source lines 10171--10171. Cross-chapter standard prerequisite located in the project tag list staged in the local snapshot; invoked at proof line 13659 of the counted proof. It is retained only as a navigation pointer; every load-bearing step is inside the delivered span.\par\medskip
\ifdefined\hypertarget\hypertarget{sourcepointer-10}{}\fi
\noindent\textit{Source reference.} Original lemma 01LH: sources/stacks-chow/tags.txt, source lines 2086--2086. Cross-chapter standard prerequisite located in the project tag list staged in the local snapshot; invoked at proof line 13676 of the counted proof. It is retained only as a navigation pointer; every load-bearing step is inside the delivered span.\par\medskip
\begin{lemma}
\label{lemma-complex-and-divisor-blowup-pre}
In Situation \ref{situation-complex-and-divisor}
there is a unique morphism $b : X' \to X$ such that
\begin{enumerate}
\item the pullback $D' = b^{-1}D$ is defined and
$(X', D', M')$ is a good triple where $M' = Lb^*M$, and
\item for any morphism of schemes $h : Y \to X$ such that
the pullback $E = h^{-1}D$ is defined and $(Y, E, Lh^*M)$
is a good triple, there is a unique factorization of $h$ through $b$.
\end{enumerate}
Moreover, for any affine chart $(U, A, f, M^\bullet)$ the restriction
$b^{-1}(U) \to U$ is the blowing up in the product of the ideals
$I_i(M^\bullet, f)$ and for any quasi-compact open $W \subset X$ the
restriction $b|_{b^{-1}(W)} : b^{-1}(W) \to W$ is a $W \setminus D$-admissible
blowing up.
\end{lemma}

\begin{proof}
The proof is just that we will locally blow up $X$ in the product ideals
$I_i(M^\bullet, f)$ for any affine chart $(U, A, f, M^\bullet)$.
The first few lemmas in More on Algebra, Section
\ref{more-algebra-section-perfect-eta} show that this is well defined.
The universal property (2) then follows from the
universal property of blowing up.
The details can be found below.

\medskip\noindent
Let $U, A, f, M^\bullet$ be an affine chart for $(X, D, M)$.
All but a finite number of the ideals $I_i(M^\bullet, f)$ are
equal to $A$ hence it makes sense to look at
$$
I = \prod\nolimits_i I_i(M^\bullet, f)
$$
and this is a finitely generated ideal of $A$. Denote
$$
b_U : U' \to U
$$
the blowing up of $U$ in $I$. Then $b_U^{-1}(U \cap D)$ is defined by
Divisors, Lemma
\ref{divisors-lemma-blow-up-pullback-effective-Cartier}.
Recall that $f^{r_i} \in I_i(M^\bullet, f)$ and hence
$b_U$ is a $(U \setminus D)$-admissible blowing up.
By Divisors, Lemma \ref{divisors-lemma-blowing-up-two-ideals}
for each $i$ the morphism $b_U$ factors as $U' \to U'_i \to U$
where $U'_i \to U$ is the blowing up in $I_i(M^\bullet, f)$
and $U' \to U'_i$ is another blowing up. It follows that
the pullback $I_i(M^\bullet, f)\mathcal{O}_{U'}$
of $I_i(M^\bullet, f)$ to $U'$ is an invertible ideal sheaf, see
Divisors, Lemmas \ref{divisors-lemma-blow-up-pullback-effective-Cartier} and
\ref{divisors-lemma-blowing-up-gives-effective-Cartier-divisor}.
It follows that $(U', b^{-1}D, Lb^*M|_U)$ is a good triple, see
Lemma \ref{lemma-pullback-triple-ideals-good}
for the behaviour of the ideals $I_i(-,-)$ under pullback.
Finally, we claim that $b_U : U' \to U$ has the universal property
mentioned in part (2) of the statement of the lemma. Namely, suppose
$h : Y \to U$ is a morphism of schemes such that
the pullback $E = h^{-1}(D \cap U)$ is defined and $(Y, E, Lh^*M)$
is a good triple. Then $Y$ is covered by affine charts
$(V, B, g, N^\bullet)$ such that $I_i(N^\bullet, g)$ is
an invertible ideal for each $i$. Then $g$ and the image of
$f$ in $B$ differ by a unit as they both cut out the effective
Cartier divisor $E \cap V$. Hence we may assume $g$ is the image of $f$ by
More on Algebra, Lemma \ref{more-algebra-lemma-eta-change-unit}.
Then $I_i(N^\bullet, g)$ is isomorphic to $I_i(M^\bullet \otimes_A B, g)$
as a $B$-module by More on Algebra, Lemma
\ref{more-algebra-lemma-ideal-well-defined}.
Thus $I_i(M^\bullet \otimes_A B, g) = I_i(M^\bullet, f)B$
(Lemma \ref{lemma-pullback-triple-ideals-good})
is an invertible $B$-module.
Hence the ideal $IB$ is invertible. It follows that $I\mathcal{O}_Y$
is invertible. Hence we obtain a unique factorization of $h$ through
$b_U$ by Divisors, Lemma \ref{divisors-lemma-universal-property-blowing-up}.

\medskip\noindent
Let $\mathcal{B}$ be the set of affine opens $U \subset X$
such that there exists an affine chart $(U, A, f, M^\bullet)$
for $(X, D, M)$. Then $\mathcal{B}$ is a basis for the topology
on $X$; details omitted. For $U \in \mathcal{B}$ we have
the morphism $b_U : U' \to U$ constructed above which satisfies
the universal property over $U$. If $U_1 \subset U_2 \subset X$
are both in $\mathcal{B}$, then $b_{U_1} : U'_1 \to U_1$
is canonically isomorphic to
$$
b_{U_2}|_{b_{U_2}^{-1}(U_1)} : b_{U_2}^{-1}(U_1) \longrightarrow U_1
$$
by the universal property. In other words, we get an isomorphism
$U'_1 \to b_{U_2}^{-1}(U_1)$ over $U_1$. These isomorphisms satisfy
the cocycle condition (again by the universal property) and hence by
Constructions, Lemma \ref{constructions-lemma-relative-glueing}
we get a morphism $b : X' \to X$ whose restriction to each
$U$ in $\mathcal{B}$ is isomorphic to $U' \to U$.
Then the morphism $b : X' \to X$ satisfies properties (1) and (2)
of the statement of the lemma as these properties may be checked
locally (details omitted).

\medskip\noindent
We still have to prove the final assertion of the lemma.
Let $W \subset X$ be a quasi-compact open.
Choose a finite covering $W = U_1 \cup \ldots \cup U_T$
such that for each $1 \leq t \leq T$ there exists an affine chart
$(U_t, A_t, f_t, M_t^\bullet)$. We will use below that for any affine
open $V = \Spec(B) \subset U_t \cap U_{t'}$ we have
(a) the images of $f_t$ and $f_{t'}$ in $B$ differ by a unit, and
(b) the complexes $M_t^\bullet \otimes_A B$ and $M_{t'} \otimes_A B$
define isomorphic objects of $D(B)$. For $i \in \mathbf{Z}$, set
$$
N_i = \max\nolimits_{t = 1, \ldots, T}
\left(\sum\nolimits_{j \geq i} (-1)^{j - i} rk(M_t^j)\right)
$$
Then $N_t - \sum\nolimits_{j \geq i} (-1)^{j - i} rk(M_t^j) \geq 0$
and we can consider the ideals
$$
I_{t, i} =
f_t^{N_i - \sum\nolimits_{j \geq i} (-1)^{j - i} rk(M_t^j)}
I_i(M_t^\bullet, f_t)
\subset A_t
$$
It follows from More on Algebra, Lemmas
\ref{more-algebra-lemma-eta-change-unit} and
\ref{more-algebra-lemma-ideal-well-defined}
that the ideals $I_{t, i}$ glue to a quasi-coherent, finite type
ideal $\mathcal{I}_i \subset \mathcal{O}_W$. Moreover, all but
a finite number of these ideals are equal to $\mathcal{O}_W$.
Clearly, the morphism $X' \to X$ constructed above
restricts to the blowing up of $W$ in the product
of the ideals $\mathcal{I}_i$. This finishes the proof.
\end{proof}

\end{document}
